\begin{abstract}
Los satélites definidos por software (SDS) representan un paradigma transformador en las comunicaciones espaciales al permitir la reconfiguración dinámica de payloads y funciones de radio a través de actualizaciones de software. Este artículo explora la integración de técnicas de Inteligencia Artificial (IA) en el procesamiento de señales digitales (DSP) para SDS, con énfasis en el aumento de flexibilidad operativa, adaptabilidad espectral y resiliencia en entornos dinámicos. Se presenta una revisión del estado del arte en IA para radios definidas por software (SDR) en contextos satelitales, una propuesta de arquitectura cognitiva basada en aprendizaje profundo para clasificación de modulación, estimación de parámetros y mitigación de interferencias en tiempo real. Los resultados de simulaciones y pruebas con datasets satelitales reales demuestran mejoras significativas en precisión de reconocimiento (hasta 15-20% sobre métodos tradicionales), eficiencia espectral y reducción de latencia en reconfiguración. Se discuten desafíos de implementación onboard, incluyendo restricciones computacionales y robustez ante efectos del canal espacial. Finalmente, se delinean direcciones futuras hacia SDS cognitivos autónomos para constelaciones LEO y misiones multi-propósito.
\end{abstract}

\begin{IEEEkeywords}
Artificial Intelligence (AI); Cognitive Satellite Architecture; Deep Learning; Digital Signal Processing (DSP); Interference Mitigation; LEO Constellations; Modulation Classification; Onboard Processing; Satellite Communications; Software-Defined Satellites (SDS).
\end{IEEEkeywords}